\section{Conclusion}
This work evaluated static prediction of flaky tests in C\# from code vocabulary. The lexical signal exists, but it is modest: across the 17 projects, which total 19,211 executed test cases, corresponding to 12,665 labeled test methods, of which 30 are flaky, the best models reach an MCC of 0.51--0.53 and an AUC of 0.86--0.87 under cross-validation, and an AUC of 0.75--0.78 when the repository is unseen, with a precision of 2--3\% at the natural base rate. The signal rests on Task-based concurrency, that is, operations run in the background and awaited with async/await, on time handling, and on resource lifecycles, and it is useful for prioritizing reruns, not for gating pipelines. Seven of the 30 flaky tests fall outside the top quartile of their repository's ranking, meaning the model does not place them among those that would be rerun first. Reading the code of some of these tests, we noticed that, in certain cases, the cause of flakiness is not in the test body, which is what the vocabulary sees: the test merely calls a helper method where the wait happens, or depends on garbage collector behavior or on state shared between tests. The model also gives high scores to stable tests that use time and concurrency in a controlled way, such as waiting for events and simulated clocks, because the words they are made of resemble those of flaky tests, even though their usage is correct. This is a qualitative analysis of a few examples, not a count of causes. As next steps, we plan to make the model follow calls between methods, so that it also reads the code of the helper methods the test uses; to design features that distinguish a fixed wait from time controlled by the test itself; to repeat the study with more projects, chosen at random; and to ask practitioners how they would use a test ranking to decide what to rerun.